\chapter{Glosarium}

\begin{description}
  \setlength{\itemsep}{2pt}
  \item[AI generatif (GenAI)] Kecerdasan artifisial yang menghasilkan konten baru, misalnya teks, gambar, atau kode.
  \item[Algoritma] Urutan langkah yang jelas untuk menyelesaikan suatu masalah.
  \item[Argumen] Nilai yang dikirim ke fungsi saat fungsi dipanggil.
  \item[\textit{Assertion}] Pernyataan dalam pengujian yang memeriksa apakah hasil sesuai dengan harapan.
  \item[Atribut] Data yang dimiliki sebuah objek atau kelas.
  \item[\textit{Boolean}] Tipe data yang hanya bernilai benar atau salah, ditulis \texttt{True} dan \texttt{False}.
  \item[\textit{Breakpoint}] Titik di dalam program tempat eksekusi dihentikan sementara agar keadaan program dapat diperiksa.
  \item[\textit{Bytecode}] Bentuk antara hasil penerjemahan kode Python yang dijalankan oleh \textit{interpreter}.
  \item[\textit{Capstone project}] Proyek penutup yang menggabungkan seluruh materi dalam satu karya.
  \item[\textit{Compiler}] Program yang menerjemahkan seluruh kode sumber menjadi berkas yang dapat dijalankan.
  \item[CSV] \textit{Comma-Separated Values}, format berkas teks yang menyimpan data tabel dengan tanda pemisah antarkolom, umumnya koma.
  \item[\textit{Debugging}] Proses menemukan, memahami, dan memperbaiki kesalahan pada program.
  \item[\textit{Dictionary}] Struktur data yang menyimpan pasangan kunci dan nilai.
  \item[Efikasi diri] Keyakinan seseorang atas kemampuannya menyelesaikan suatu tugas.
  \item[\textit{Exception}] Kesalahan yang muncul saat program berjalan dan dapat ditangani oleh program.
  \item[\textit{F-string}] Penulisan teks berawalan \texttt{f} yang dapat memuat nilai variabel di dalam kurung kurawal.
  \item[\textit{Float}] Tipe data bilangan desimal.
  \item[Fungsi] Blok kode bernama yang dapat dipanggil berulang kali.
  \item[GUI] \textit{Graphical User Interface}, antarmuka program berupa jendela, tombol, dan komponen grafis lain.
  \item[IDE] \textit{Integrated Development Environment}, perangkat lunak yang menyatukan editor kode, alat untuk menjalankan program, dan alat \textit{debugging}.
  \item[Indentasi] Jarak kosong di awal baris yang menandai blok kode di Python.
  \item[Indeks] Nomor posisi elemen dalam sebuah urutan, dimulai dari 0.
  \item[\textit{Integer}] Tipe data bilangan bulat.
  \item[\textit{Interpreter}] Program yang membaca dan menjalankan kode secara langsung.
  \item[Iterasi] Satu kali putaran dalam perulangan.
  \item[Kelas] Rancangan yang mendefinisikan atribut dan metode bagi objek.
  \item[Konstanta] Variabel yang nilainya disepakati tidak berubah selama program berjalan.
  \item[Konstruk] Variabel yang tidak dapat diamati langsung sehingga diukur lewat beberapa indikator, misalnya butir kuesioner.
  \item[Konversi tipe data] Pengubahan nilai dari satu tipe data ke tipe data lain.
  \item[Kuasi-eksperimen] Rancangan penelitian yang membandingkan kelompok perlakuan dan kelompok pembanding tanpa mengacak peserta satu per satu.
  \item[\textit{Large language model} (LLM)] Model AI yang dilatih dengan teks dalam jumlah sangat besar sehingga dapat membuat dan menafsirkan teks, termasuk kode program.
  \item[\textit{List}] Struktur data berurutan yang isinya dapat diubah.
  \item[\textit{Logic error}] Kesalahan yang membuat program berjalan tanpa pesan kesalahan, tetapi hasilnya salah.
  \item[Metode] Fungsi yang dimiliki sebuah kelas.
  \item[Modul] Berkas Python berisi kode yang dapat diimpor oleh program lain.
  \item[Objek] Wujud nyata dari sebuah kelas yang menyimpan datanya sendiri.
  \item[Operator] Simbol atau kata kunci untuk melakukan operasi terhadap nilai.
  \item[Paket] Folder berisi modul yang dapat diimpor sebagai satu kesatuan.
  \item[Parameter] Nama pada definisi fungsi yang menampung nilai dari pemanggil.
  \item[Percabangan] Struktur program yang memilih baris yang dijalankan berdasarkan kondisi.
  \item[Perulangan] Struktur program yang menjalankan blok kode berkali-kali.
  \item[Pewarisan] Mekanisme kelas turunan memperoleh atribut dan metode dari kelas induk.
  \item[PLS-SEM] \textit{Partial least squares structural equation modeling}, teknik analisis untuk menguji hubungan antarkonstruk sekaligus menilai indikator pengukurnya.
  \item[Polimorfisme] Kemampuan objek yang berbeda merespons pemanggilan metode yang sama dengan caranya masing-masing.
  \item[\textit{Prompt}] Perintah atau pertanyaan tertulis yang diberikan kepada model AI.
  \item[\textit{Prompt} terarah] \textit{Prompt} yang meminta AI memberi petunjuk dan memeriksa jawaban tanpa menuliskan penyelesaiannya.
  \item[\textit{Range}] Objek yang menghasilkan deret angka.
  \item[\textit{Refactoring}] Perubahan struktur kode tanpa mengubah perilakunya.
  \item[Rekursi] Teknik ketika sebuah fungsi memanggil dirinya sendiri.
  \item[\textit{Runtime error}] Kesalahan yang baru muncul saat program berjalan.
  \item[\textit{Set}] Kumpulan elemen unik tanpa urutan.
  \item[\textit{String}] Tipe data teks.
  \item[\textit{Syntax error}] Kesalahan penulisan kode yang terdeteksi sebelum program berjalan.
  \item[\textit{Traceback}] Laporan urutan pemanggilan yang ditampilkan saat terjadi kesalahan.
  \item[\textit{Tuple}] Struktur data berurutan yang isinya tidak dapat diubah.
  \item[\textit{Unit test}] Pengujian terhadap bagian kecil program secara terpisah.
  \item[Variabel] Nama yang menyimpan sebuah nilai selama program berjalan.
  \item[\textit{Widget}] Komponen antarmuka grafis, misalnya tombol atau kotak isian.
\end{description}
